\section{Introducción: por qué K2 es el «cerebro en una cubeta» al que le crecen manos}

Esta sección establece primero el hilo técnico de todo el episodio. Kimi K2 es un modelo de lenguaje de gran escala de código abierto, orientado a la programación y a lo Agentic, basado en la arquitectura MoE. La metáfora más contundente de la entrevista es la del «cerebro en una cubeta»: un modelo que solo sabe pensar y producir texto es como un cerebro colocado en una pecera, sin una conexión real con el mundo exterior; en cambio, cuando el modelo adquiere capacidad de programación, capacidad de usar herramientas en múltiples turnos y retroalimentación externa, empiezan a crecerle «manos» y puede manipular el mundo digital.

En este episodio hay que entender cuatro cosas. Primero, por qué K2 reúne el modelo base, la token efficiency, las capacidades Agentic y el código abierto. Segundo, por qué la test-time scaling apunta a la vez al RL de pensamiento largo y al Agent RL. Tercero, cómo se equilibran el muro de datos, el FLOPs scaling, la Linear Attention, las arquitecturas de contexto largo y el «coeficiente intelectual» del modelo. Cuarto, cómo Yang Zhilin (杨植麟) integra el juicio técnico, el modelo de negocio y la mentalidad de fundador en el relato de «estar al comienzo del infinito».

\begin{figure}[H]
\centering
\includegraphics[width=0.96\textwidth]{brain-in-vat-to-agent.png}
\caption{Del cerebro en una cubeta al Agent: la capacidad de programación y el uso de herramientas permiten que el modelo empiece a manipular el mundo digital. Diagrama conceptual propio, elaborado a partir de la conversación de 00:01:19--00:24:58.}
\end{figure}

\begin{knowledgebox}{Lectura de la figura: por qué las «manos» son decisivas}
Que un modelo sepa pensar y que sepa actuar son dos cosas distintas. El código, las llamadas a herramientas, las API, el sistema de archivos y los entornos de retroalimentación permiten que el razonamiento del modelo modifique el estado externo. Lo esencial de un Agentic LLM no es solo decir más pasos, sino usar herramientas en múltiples turnos y corregirse según los resultados.
\end{knowledgebox}

\begin{importantbox}{Tesis central del episodio}
La importancia de K2 no es solo ser «otro modelo potente», sino integrar en un mismo sistema el modelo base, la capacidad de programación, el uso de herramientas, el ecosistema de código abierto y la test-time scaling: llevar al modelo de un sistema cognitivo cerrado a un agente digital capaz de actuar.
\end{importantbox}

\begin{knowledgebox}{Panorama del sistema en este episodio}
\begin{tabularx}{\textwidth}{p{0.20\textwidth}p{0.34\textwidth}X}
\toprule
Nivel & Pregunta central & Respuesta en la entrevista \\
\midrule
Nivel filosófico & ¿Por qué seguir escalando? & Los problemas son inevitables, pero los problemas pueden resolverse. \\
Nivel del modelo & ¿Cómo se convierte K2 en una base sólida? & MoE, token efficiency, Muon, Rephrase. \\
Nivel de acción & ¿Cómo le crecen manos al modelo? & Capacidad de programación, uso de herramientas en múltiples turnos, generalización Agentic. \\
Nivel del ecosistema & ¿Por qué código abierto? & Para que las capacidades del modelo base se validen externamente y el ecosistema las amplifique. \\
Nivel comercial & ¿Cómo se gana dinero y dónde están los límites? & API, productos Agent y una nueva delimitación entre empresas de modelos base y empresas de aplicaciones. \\
Nivel organizacional & ¿Cómo soporta el fundador la complejidad? & Gestionar con retroalimentación al estilo RL, pero con cuidado de no caer en el hack de las métricas. \\
\bottomrule
\end{tabularx}
\end{knowledgebox}

\begin{knowledgebox}{Términos clave: palabras clave del episodio}
\begin{tabularx}{\textwidth}{p{0.22\textwidth}p{0.32\textwidth}X}
\toprule
Término & Significado & Por qué importa \\
\midrule
MoE & Mixture of Experts: se enruta cada token o tarea a distintos expertos & Aumenta la capacidad de parámetros y, al mismo tiempo, controla el cómputo activado en cada paso. \\
Agentic LLM & Modelo de lenguaje que puede usar herramientas en múltiples turnos, observar la retroalimentación y completar tareas & Pasa de la conversación a la acción en el mundo digital. \\
Test-time scaling & Invertir más cómputo, más pensamiento o más llamadas a herramientas durante la inferencia & Permite que el modelo siga ampliando sus capacidades en la fase de prueba. \\
Token efficiency & La ganancia de capacidad que aporta cada token de entrenamiento & En la era del muro de datos, con los mismos datos el «cerebro» debe crecer más. \\
Muon & Línea de optimizadores mencionada en la entrevista & Puede mejorar la eficiencia del entrenamiento, pero plantea retos de estabilidad. \\
\bottomrule
\end{tabularx}
\end{knowledgebox}

\subsection{Resumen de la sección}

EP113 es una entrevista de juicio técnico posterior al lanzamiento de K2. Sitúa al Agentic LLM dentro de la estrategia de las empresas de modelos, el código abierto, la eficiencia del entrenamiento, las decisiones de arquitectura y la psicología del fundador, y conviene leerla junto con la revisión de arquitecturas de Attention de EP119 y el marco teórico de Agent de EP115.
